\chapter*{परिचय} \label{chap:pref}
\addcontentsline{toc}{chapter}{परिचय}
ब्रह्मांड ३ गुणों से बना है: \bt{\satva{सत्व}, \rajas{रजस्}, \tamas{तमस्}}। हर आत्मा के मन पर यह गुण प्रभाव डालकर उसमे भिन्न प्रकार की इच्छाएँ पैदा करती हैं, जिनकी प्राप्ति हेतु आत्मा अपने शरीर के द्वारा कर्म करके उन कर्मों से उत्पन्न अनुभव को भुगतती है। कर्मों के फलित होने की यात्रा हमारे बुद्धि द्वारा समझना  उतना ही सरल है जितना इस जगत् का वास्तविक सत्य सिर्फ विज्ञान से समझना होगा - बिल्कुल नहीं। किसीको पता नहीं कि हमारे सैकडो जन्मो के कर्मों मे से कौनसा किस रूप मे और कैसे हमे वापस मिलने आएगा, हम बस भगवद्गीता के द्वारा इतना जान सकते हैं कि किस कर्म का कैसा परिणाम मिलेगा।

जब तक एक आत्मा खुदको परमात्मा का हिस्सा नहीं जान लेती और अपने पृथ्वी पर शेष कर्म खत्म नहीं करती, तब तक उसका पुनर्जन्म निश्चिततः होते रहेगा। इस यात्रा पर  उसे बहुत कष्ट, पीड़ा और अस्थायी खुशियों का सामना करना पड़ता हैं, और बहुत बार, तामसिक और राजसिक गुणों के गहरे चपेट मे आए हुए के लिए स्थिती काफी गम्भीर भी हो सकती है। आत्मा वास्तव में होती तो अमर ही है, और हर एक को एक दिन इस सत्य की पहचान निश्चिततः होगी, परन्तु इस यात्रा पर एक अनभिज्ञ तभी लग पाएगा जब उसे कोई मार्गदर्शन मिले। इस मार्गदर्शन  के बहुत स्वरूप है, जिन्हें गीता में अलग-अलग योगमार्ग के द्वारा समझाया गया है।

कुछ लोग समर्पण और भक्ति चुनते हैं, तो कुछ ज्ञान में अपना उत्तर पाते हैं। कुछ आत्मनियंत्रण और अनासक्त कर्म करना पसंद करते हैं तो कुछ सन्यासी जैसे खुदको दुनियादारी से सीधे निकाल ही लेते हैं। सब जा वही रहे हैं, परन्तु अलग व्यक्तियों के स्वभाव के हिसाब से उन्हें अलग मार्ग ज्यादा स्वाभाविक और सहज लगेंगे। कोई मार्ग गलत नहीं है, और न ही किसीको एक अस्वाभाविक मार्ग का पालन करने विवश करना सही है।

जो वैज्ञानिक मन वाली आत्माएँ अध्यात्म और सत्य के खोज में पहली बार प्रवेश करती हैं, उनके सामने एक अनन्त गहराई का समन्दर प्रकट होगा। हजारो विचारधाराएँ, सैकडो मान्यताएँ, और उतने ही विरोधाभास, जहाँ धर्म की हमेशा जीत होगी यह बोलने, मानने और उस हिसाब व्यवहार करने वालों के साथ ही कई बार कठोर  अन्याय होता है, जबकी अधर्मी और पापी लोग कानून और न्याय के ऊपर बैठे ऐसे प्रतीत होते हैं, जैसे की  वह भगवान ही बने हो, और जिनके लिए सामान्य लोग बस खिलौने का मोल रखते हो।

इन विरोधाभासों का हल न मिले तब तक आजके संसार में एक वैज्ञानिक के लिए धर्म तथा ब्रह्मांड की गुप्त कर्मप्रणाली के न्याय पर विश्वास रखना असम्भव है। गीता के द्वारा ज्ञान पाना कि कर्मों के फल हमेशा मिलेंगे इसका एक स्थान है, परन्तु इन कर्मों से भाग्य कैसे लिखा जाता है, भुगतना पड़ता है, तथा हमारा भाग्य अपने कर्मों से कैसे बदला जा सकता है इसको प्रत्यक्ष देख पाने का अपना अनोखा स्थान है। ज्योतिष्शास्त्र एक ऐसी प्रणाली है, जिससे इसी कर्मचक्र का प्रत्यक्ष दर्शन किया  जा सकता है।

इस पुस्तक के दो खण्ड है -
\begin{enumerate}
\item \autoref{part:calc} में आकाश में स्थित ज्योतियों (ग्रह, नक्षत्र, राशी इत्यादि) कि वैज्ञानिक तौर पर चर्चा होगी।  ग्रहों की पृथ्वी से दिखने वाली गति, ज्योतियों के प्रकार तथा पंचांग से कौनसा ग्रह आकाश में कब कहाँ और कैसे प्रतीत होगा इसका गणित इस अध्याय का भाग है। सारी कुन्डलियाँ इसी गणित के आधार पर बनाई जाती है। आजकल यह गणित  सॉफ्टवेअर के द्वारा किया जाता है, परन्तु  इस गणना के बिना \autoref{part:pred} मे दी गई सारी विधियाँ बेकार होंगी। यह अध्याय ज्योतिष्शास्त्र का मूलाधार है।
\item \autoref{part:pred} में कुन्डलियों के प्रकार, उनकी क्षमताएँ तथा उनसे की जा सकने वाली हर तरह की भविष्यवाणियों और अनुमानों का वर्णन है। पहले अध्याय से अलग,  यह अध्याय औपचारिक रूप से अनन्त है, क्योंकि करोडो सूत्र लिखे जा चुके है, परन्तु यहाँ बस वही विधियाँ लिखी गई हैं, जो आजके कलियुग में सटीक रूप से काम आए, और मुझे पता हो। 
\end{enumerate}
यह गुप्त शास्त्र अध्यात्मिक को ही कभी काम आ सकता है, और यदि उसके उपयोगी सूत्र इस पुस्तक मे एकजुत करने से एक  भी आत्मा को  सही कर्म करने की प्रेरणा मिलती है, तो उसका ध्येय सफल हो चुका होगा।